% Part: computability
% Chapter: computability-theory
% Section: equiv-ce-defs

\documentclass[../../../include/open-logic-section]{subfiles}

\begin{document}

\olfileid[fr]{cmp}{thy}{eqc}
\olsection[Définitions des ensembles effectivement énumérables]{Caractérisations équivalentes des ensembles effectivement énumérables}

Les énoncés suivants donnent plusieurs façons
importantes de caractériser les ensembles effectivement
énumérables.

\begin{thm}
\ollabel{thm:ce-equiv}
Soit $S$ un ensemble d'entiers naturels. Les propriétés
suivantes sont équivalentes :
\begin{enumerate}
\item\ollabel{case:ce} $S$ est effectivement énumérable.
\item\ollabel{case:ran-pc} $S$ est l'image d'une fonction
  calculable \emph{partielle}.
\item\ollabel{case:ran-prim} $S$ est vide ou l'image
  d'une fonction récursive primitive.
\item\ollabel{case:ce-domain} $S$ est le \emph{domaine}
  d'une fonction calculable partielle.
\end{enumerate}
\end{thm}

\begin{explain}
Les trois premières propriétés disent qu'un ensemble
effectivement énumérable non vide peut être énuméré
aussi bien par une fonction calculable que par une
fonction calculable partielle ou une fonction récursive
primitive. La quatrième nous dit que, si $S$ est
effectivement énumérable, il existe un indice~$e$ tel que
\[
S = \Setabs{x}{\cfind{e}(x) \fdefined}.
\]
Autrement dit, $S$ est l'ensemble des entrées sur
lesquelles le calcul de $\cfind{e}$ s'arrête. C'est
pourquoi on appelle aussi ces ensembles
\emph{semi-décidables} : lorsqu'un nombre appartient à
l'ensemble, on finit par obtenir « oui »; sinon,
la procédure ne donne jamais « non ».
\end{explain}

\begin{proof}
Toute fonction récursive primitive est calculable,
et toute fonction calculable est, en particulier, une fonction partielle calculable.
Ainsi, \olref{case:ran-prim} implique \olref{case:ce},
qui implique à son tour \olref{case:ran-pc}.
Si $S$ est vide, c'est l'image de la fonction calculable
partielle qui n'est définie nulle part. Il suffit donc
de montrer que \olref{case:ran-pc} implique
\olref{case:ran-prim} pour établir l'équivalence
des trois premières propriétés.

Supposons que $S$ soit l'image de la fonction calculable
partielle $\cfind{e}$. Si $S$ est vide, c'est fini.
Sinon, choisissons $a\in S$. Par le théorème de forme
normale de Kleene, on peut écrire
\[
\cfind{e}(x) \simeq U(\umin{s}{T(e, x, s)}).
\]
En particulier, $\cfind{e}(x) \fdefined$ et vaut~$y$
si et seulement s'il existe un~$s$ tel que $T(e,x,s)$
et $U(s)=y$. Définissons $f(z)$ par
\[
f(z) = \begin{cases}
  U((z)_1) & \text{si $T(e, (z)_0, (z)_1)$} \\
  a        & \text{sinon.}
\end{cases}
\]
La fonction $f$ est récursive primitive, puisque
$T$ et~$U$ le sont. \iftag{TMs}{En termes de machines
  de Turing, on lit dans $z$ les valeurs $(z)_0$ et $(z)_1$.
  Si $(z)_1$ décrit
  un calcul qui s'arrête pour la machine~$M_e$ sur
  l'entrée $(z)_0$, alors $f$ renvoie la sortie de ce
  calcul; sinon, elle renvoie~$a$.}{}

Montrons que l'image de $f$ est bien~$S$. Si $y\in S$,
il existe $x$ et~$s$ tels que $T(e,x,s)$ et $U(s)=y$;
on a donc $y=f(\tuple{x,s})$. Réciproquement, si $y$
est dans l'image de $f$, alors soit $y=a$, soit il
existe un~$z$ tel que $T(e,(z)_0,(z)_1)$ et
$U((z)_1)=y$. Dans ce dernier cas,
$\cfind{e}((z)_0)\fdefined=y$; dans les deux cas,
$y\in S$.

(La notation $\cfind{e}(x) \fdefined = y$ signifie
« $\cfind{e}(x)$ est définie et vaut $y$ ». On
pourrait aussi écrire $\cfind{e}(x)=y$, mais la
flèche rappelle utilement qu'il s'agit d'une fonction
partielle.)\footnote{Corrections éditoriales : dans le sens réciproque du raisonnement sur l'image de $f$, la source emploie un $x$ non fixé là où il faut $(z)_0$. Dans le passage de l'énumération au domaine, nous traitons aussi explicitement l'ensemble vide, qui ne peut être l'image d'une fonction totale sur $\Nat$. Nous écrivons $\simeq$ dans la forme normale pour couvrir les valeurs indéfinies.}

Pour achever la preuve du \olref{thm:ce-equiv},
montrons que \olref{case:ce} et
\olref{case:ce-domain} sont équivalents. Supposons
d'abord \olref{case:ce}. Si $S$ est vide, la
fonction calculable partielle nulle part définie a
pour domaine~$S$. Sinon, $S$ est l'image d'une
fonction calculable~$f$, c'est-à-dire
\[
S = \Setabs{y}{\text{il existe $x$ tel que } f(x) = y}.
\]
Posons
\[
g(y) = \umin{x}{(f(x) = y)}.
\]
La fonction $g$ est calculable partielle, et $g(y)$
est définie si et seulement s'il existe un~$x$ tel
que $f(x)=y$. Son domaine est donc l'image de~$f$.
\iftag{TMs}{En termes de machines de Turing, si une
machine~$F$ énumère les éléments de~$S$, une machine~$G$
semi-décide l'appartenance à~$S$ en cherchant un
élément donné parmi les sorties de~$F$. Elle s'arrête
si elle le trouve et poursuit la recherche indéfiniment
sinon.}{}

Enfin, supposons \olref{case:ce-domain} : $S$ est
le domaine d'une fonction calculable partielle
$\cfind{e}$, soit
\[
S = \Setabs{x}{\cfind{e}(x) \fdefined}.
\]
Si $S$ est vide, le résultat est immédiat. Sinon,
choisissons $a\in S$ et définissons $f$ par
\[
f(z) = \begin{cases}
(z)_0 & \text{si $T(e,(z)_0,(z)_1)$} \\
a & \text{sinon.}
\end{cases}
\]
Comme précédemment, un entier~$x$ appartient à
l'image de~$f$ si et seulement si
$\cfind{e}(x)\fdefined$, c'est-à-dire si et
seulement si $x\in S$. \iftag{TMs}{En termes de
machines de Turing, étant donnée une machine~$M_e$
qui semi-décide~$S$, on énumère les éléments de~$S$
en parcourant tous les calculs possibles et en
renvoyant les entrées des calculs qui s'arrêtent.}{}
\end{proof}

La propriété~\olref{case:ce-domain} du
\olref{thm:ce-equiv} permet d'énumérer commodément
les ensembles effectivement énumérables : pour chaque
$e$, notons $W_e$ le domaine de $\cfind{e}$, soit
\[
W_e = \Setabs{x}{\cfind{e}(x) \fdefined}.
\]
Tout ensemble effectivement énumérable~$A$ est alors
égal à $W_e$ pour un certain~$e$.

Le théorème suivant en donne encore une autre
caractérisation.

\begin{thm}
\ollabel{thm:exists-char}
Un ensemble~$S$ est effectivement énumérable si et
seulement s'il existe une relation calculable~$R(x,y)$
telle que
\[
S = \Setabs{ x }{ \lexists[y][R(x,y)] }.
\]
\end{thm}

\begin{proof}
Dans un sens, si $S$ est effectivement énumérable,
il existe un~$e$ tel que $S=W_e$. Pour cet indice,
on peut écrire
\[
S = \Setabs{ x }{ \lexists[y][T(e, x, y)] }.
\]
Dans l'autre sens, supposons
$S=\Setabs{ x }{ \lexists[y][R(x, y)] }$.
Définissons $f$ par
\[
f(x) \simeq \umin{y}{R(x, y)}.
\]
Alors $f$ est calculable partielle et son domaine
est~$S$.
\end{proof}

\end{document}
